%! Simplifying Rational Expressions — Unit 5, Lesson 5.1 — Exit Ticket
\documentclass[10pt]{article}
\usepackage{algebra2-article}
\usepackage{algebra2-boxes}
\newcommand{\TallMath}[1]{$\displaystyle #1\rule[-1.4em]{0pt}{3.2em}$}

\begin{document}

\pageheader{Unit 5, Lesson 5.1}{Exit Ticket}
\namedateperiod

\vspace{0.08in}

No notes. Show the factored line for question 1.

\vspace{0.06in}

\begin{enumerate}[leftmargin=1.8em, itemsep=9pt]
  \item \textbf{Simplify and restrict.} \; $\dfrac{x^2-2x-15}{x^2-9}$

        \vspace{4pt}
        Factored form: \blank{5.0cm}

        \vspace{4pt}
        Simplest form: \blank{2.6cm} \quad Restrictions: $x\neq$ \blank{2.4cm}

        \vspace{4pt}
        One restriction is a \emph{hole} and one is a \emph{vertical asymptote}. \\[3pt]
        Hole at $x=$ \blank{1.2cm} \quad Vertical asymptote at $x=$ \blank{1.2cm}

  \item \textbf{Explain.} A classmate says: ``Once the $(x+3)$ cancels, the expression is just
        $\dfrac{x-5}{x-3}$, so $x=-3$ is fine now.'' Is the classmate right? Explain what cancelling
        does and does not change.
        \writelines{3}

  \item \textbf{SOL-style multiple choice.} Which is the correct simplest form of
        $\dfrac{4-x^2}{x^2-x-2}$, with its restrictions?
        \begin{enumerate}[label=\Alph*., leftmargin=1.9em, itemsep=1pt]
          \item $\dfrac{x+2}{x+1}$, \; $x\neq2,\,-1$
          \item $-\dfrac{x+2}{x+1}$, \; $x\neq2,\,-1$
          \item $-\dfrac{x+2}{x+1}$, \; $x\neq-1$
          \item $-\dfrac{4}{x+2}$, \; $x\neq-2$
        \end{enumerate}
\end{enumerate}

\end{document}
